\part{Appendices}

\chapter{Glossary}

\begin{description}[leftmargin=0pt,style=nextline,font=\bfseries\color{accent}]
\item[Address space] The complete set of addresses a process can name. Each user process has its own,
  defined by its page directory.
\item[ATA / IDE] The interface used to talk to the disk. \emph{PIO} means the processor moves the
  data itself, one word at a time.
\item[BIOS] The PC firmware. Initialises hardware, loads the boot sector, and offers services that
  stop working once you enter protected mode.
\item[Bitmap (memory)] One bit per 4 KB frame recording whether it is in use.
\item[Block] The filesystem's allocation unit. 512 bytes here, the same as a disk sector.
\item[Boot sector] The first 512 bytes of a disk, ending in \hexa{AA55}. The BIOS loads it at
  \hexa{7C00}.
\item[Context] Everything that must be saved to suspend a process: registers, \reg{EIP},
  \reg{EFLAGS}, stack, address space.
\item[Context switch] Replacing the running process's context with another's.
\item[CR0 / CR2 / CR3] Control registers: mode switches, faulting address, active page directory.
\item[Descriptor] An 8-byte record in the GDT or IDT describing a segment or an interrupt handler.
\item[Directory (filesystem)] A file whose contents are (inode, name) records.
\item[DPL] Descriptor Privilege Level. The ring required to use a segment or gate.
\item[EOI] End Of Interrupt. The acknowledgement the PIC needs before it will deliver another
  interrupt.
\item[Exception] An interrupt generated by the processor itself when an instruction cannot complete.
\item[File descriptor] A small integer indexing the kernel's open-file table. 0, 1 and 2 are
  reserved for keyboard and screen.
\item[Frame] A 4 KB chunk of physical RAM.
\item[GDT] Global Descriptor Table. Defines the memory segments and their privilege levels.
\item[Heap] A region from which variable-sized allocations are made.
\item[IDT] Interrupt Descriptor Table. 256 entries mapping interrupt numbers to handlers.
\item[Identity mapping] A mapping where virtual address $X$ corresponds to physical address $X$.
\item[Inode] The structure describing a file: size, type and block pointers. Not its name.
\item[Interrupt] A mechanism by which hardware or software diverts the processor to a kernel
  routine.
\item[IRQ] A hardware interrupt line. IRQ0 is the timer, IRQ1 the keyboard.
\item[Kernel] The privileged core of the operating system, running in ring 0.
\item[LBA] Logical Block Addressing. Sectors numbered linearly from 0.
\item[Little-endian] Storing the least significant byte at the lowest address. x86 does this.
\item[Page] A 4 KB chunk of virtual address space.
\item[Page directory] A table of 1024 entries, each pointing at a page table. \reg{CR3} points at
  it.
\item[Page fault] Exception 14, raised when a translation cannot be completed.
\item[Page table] A table of 1024 entries, each describing one 4 KB page.
\item[Paging] Translating virtual addresses to physical ones in 4 KB units, enforced by hardware.
\item[PCB] Process Control Block. The structure holding everything about a process.
\item[PIC] Programmable Interrupt Controller (8259). Routes device interrupts to the CPU.
\item[PIT] Programmable Interval Timer (8253/8254). Generates periodic interrupts.
\item[PID] Process identifier.
\item[Polling] Repeatedly asking a device whether it is ready, instead of waiting for an interrupt.
\item[Preemption] Taking the processor from a process without its cooperation.
\item[Protected mode] The 32-bit x86 mode with rings, segment permissions and paging.
\item[Quantum] The number of timer ticks a process runs before being preempted.
\item[Real mode] The 16-bit x86 mode the processor starts in. 1 MB, no protection.
\item[Ring] A privilege level. 0 is the kernel, 3 is user space.
\item[Round robin] A scheduling policy giving every ready process a turn in order.
\item[RPL] Requested Privilege Level. The low two bits of a segment selector.
\item[Scancode] The number a keyboard sends for a physical key. Not a character.
\item[Scheduler] The code that decides which process runs next.
\item[Sector] The disk's physical unit, 512 bytes.
\item[Segment] A region of memory described by a GDT descriptor. Flat in this system.
\item[Selector] A 16-bit value indexing the GDT and carrying an RPL.
\item[Superblock] Block 0 of the filesystem, describing the filesystem itself.
\item[System call] A request from a user program to the kernel, made with \cod{int 0x80}.
\item[TLB] Translation Lookaside Buffer. Caches recent address translations.
\item[Triple fault] A fault while handling a double fault. The processor resets.
\item[TSS] Task State Segment. In practice, only \cod{esp0} and \cod{ss0} matter.
\item[User space] Ring 3. Where unprivileged programs run.
\item[Virtual address] An address as a program sees it, before translation.
\end{description}

% ============================================================
\chapter{x86 quick reference}

\section{General-purpose registers}

\begin{center}
\begin{tabularx}{\textwidth}{@{}l l l X@{}}
\toprule
\textbf{32-bit} & \textbf{16-bit} & \textbf{8-bit} & \textbf{Conventional use} \\
\midrule
EAX & AX & AH, AL & Return value; system call number \\
EBX & BX & BH, BL & System call argument 1 \\
ECX & CX & CH, CL & Loop counter; argument 2 \\
EDX & DX & DH, DL & I/O port; argument 3 \\
ESI & SI & --- & Source pointer \\
EDI & DI & --- & Destination pointer \\
EBP & BP & --- & Frame pointer \\
ESP & SP & --- & Stack pointer \\
\bottomrule
\end{tabularx}
\end{center}

\section{Instructions used in this project}

\begin{center}
\begin{tabularx}{\textwidth}{@{}l X@{}}
\toprule
\textbf{Instruction} & \textbf{Effect} \\
\midrule
\cod{mov d, s} & Copy \cod{s} into \cod{d} \\
\cod{push v} / \cod{pop d} & Stack: subtract 4 and store / load and add 4 \\
\cod{pusha} / \cod{popa} & Push or pop all eight general-purpose registers \\
\cod{call a} / \cod{ret} & Push return address and jump / pop and jump \\
\cod{iret} & Return from interrupt: pop EIP, CS, EFLAGS (and SS, ESP if the ring changes) \\
\cod{jmp s:o} & Far jump: change CS and EIP together \\
\cod{cmp a, b} & Subtract and discard, keeping only the flags \\
\cod{test a, b} & AND and discard, keeping only the flags \\
\cod{in d, dx} / \cod{out dx, s} & Read from / write to an I/O port \\
\cod{cli} / \cod{sti} & Clear / set the interrupt flag \\
\cod{lgdt} / \cod{lidt} / \cod{ltr} & Load GDT / IDT / task register \\
\cod{invlpg [a]} & Invalidate one TLB entry \\
\cod{lodsb} & AL = [DS:SI]; SI++ \\
\cod{hlt} & Halt until the next interrupt \\
\bottomrule
\end{tabularx}
\end{center}

\section{CPU exceptions}

\begin{center}
\begin{tabularx}{\textwidth}{@{}l l l X@{}}
\toprule
\textbf{No.} & \textbf{Name} & \textbf{Err?} & \textbf{Usual cause here} \\
\midrule
0  & Divide By Zero & no & Division by a zero variable \\
6  & Invalid Opcode & no & Jumped into data \\
8  & Double Fault & yes & A fault while handling a fault \\
13 & General Protection & yes & Privilege violation, bad selector \\
14 & Page Fault & yes & Unmapped address, or permission denied \\
\bottomrule
\end{tabularx}
\end{center}

\section{Page table entry flags}

\begin{center}
\begin{tabularx}{0.9\textwidth}{@{}l l X@{}}
\toprule
\textbf{Bit} & \textbf{Value} & \textbf{Meaning} \\
\midrule
0 & \hexa{1} & Present \\
1 & \hexa{2} & Writable \\
2 & \hexa{4} & User-accessible (ring 3) \\
5 & \hexa{20} & Accessed \\
6 & \hexa{40} & Dirty \\
\bottomrule
\end{tabularx}
\end{center}

% ============================================================
\chapter{Memory map of lytlnyblOS}

\section{Physical memory}

\begin{center}
\begin{tabularx}{\textwidth}{@{}l l X@{}}
\toprule
\textbf{Address} & \textbf{Size} & \textbf{Contents} \\
\midrule
\hexa{00000} & 1 KB & Interrupt vector table (real mode) \\
\hexa{00400} & 256 B & BIOS data area \\
\hexa{04FFC} & 2 B & \textbf{E820 entry count} (written by the bootloader) \\
\hexa{05000} & varies & \textbf{E820 memory map} \\
\hexa{07C00} & 512 B & Boot sector \\
\hexa{09000} & $\sim$23 KB & \textbf{The kernel} (\cod{\_kernel\_start}) \\
\hexa{13468} & --- & \cod{\_kernel\_end} (exact value varies per build) \\
\hexa{A0000} & 128 KB & Video memory \\
\hexa{B8000} & 4 KB & \textbf{VGA text buffer} \\
\hexa{C0000} & 256 KB & BIOS and option ROMs \\
\hexa{100000} & 4 KB & \textbf{Physical memory bitmap} \\
\hexa{101000} & rest & Free RAM for allocation \\
\bottomrule
\end{tabularx}
\end{center}

\section{Virtual memory (per process)}

\begin{center}
\begin{tabularx}{\textwidth}{@{}l l X@{}}
\toprule
\textbf{Address} & \textbf{Constant} & \textbf{Contents} \\
\midrule
\hexa{00000000} & --- & Identity-mapped kernel, first 4 MB \\
\hexa{00400000} & \cod{USER\_CODE\_BASE} & User program code \\
\hexa{00A00000} & \cod{USER\_HEAP\_START} & User heap, grows up \\
\hexa{00B00000} & \cod{USER\_VGA} & Reserved for direct user VGA access \\
\hexa{BFFFF000} & \cod{USER\_STACK\_TOP} & Top of the user stack, grows down \\
\hexa{C0000000} & \cod{HEAP\_START} & Kernel heap, grows up \\
\bottomrule
\end{tabularx}
\end{center}

\section{Disk image}

\begin{center}
\begin{tabularx}{\textwidth}{@{}l l X@{}}
\toprule
\textbf{Sector} & \textbf{Constant} & \textbf{Contents} \\
\midrule
0 & --- & Boot sector \\
1--60 & --- & Kernel (60 sectors are read by the BIOS) \\
70 & \cod{FS\_START\_BLOCK} & Filesystem superblock \\
71 & \cod{FS\_BITMAP\_BLOCK} & Block bitmap \\
72--83 & \cod{FS\_INODE\_START} & Inode table (128 inodes) \\
84+ & --- & Data blocks \\
\bottomrule
\end{tabularx}
\end{center}

% ============================================================
\chapter{System call reference}

\begin{center}
\begin{tabularx}{\textwidth}{@{}l l l X@{}}
\toprule
\textbf{EAX} & \textbf{Name} & \textbf{EBX / ECX / EDX} & \textbf{Returns} \\
\midrule
\hexa{01} & \cod{exit}   & --- & Never \\
\hexa{02} & \cod{getpid} & --- & The pid \\
\hexa{03} & \cod{yield}  & --- & --- \\
\hexa{04} & \cod{sleep}  & ticks & --- \\
\hexa{05} & \cod{write}  & fd / buffer / count & Bytes written, or $-1$ \\
\hexa{06} & \cod{read}   & fd / buffer / count & Bytes read, 0 at EOF, or $-1$ \\
\hexa{07} & \cod{sbrk}   & increment & The new heap end \\
\hexa{08} & \cod{fsops}  & operation / path & --- \\
\hexa{09} & \cod{open}   & path & A descriptor ($\geq 3$), or $-1$ \\
\hexa{0A} & \cod{close}  & fd & --- \\
\bottomrule
\end{tabularx}
\end{center}

\section{\cod{fsops} sub-operations}

\begin{center}
\begin{tabularx}{0.9\textwidth}{@{}l l X@{}}
\toprule
\textbf{EBX} & \textbf{Name} & \textbf{Effect} \\
\midrule
0 & \cod{FS\_LS} & List a directory \\
1 & \cod{FS\_MKDIR} & Create a directory \\
2 & \cod{FS\_TOUCH} & Create an empty file \\
3 & \cod{FS\_RM} & Delete a file \\
4 & \cod{FS\_RUN} & Load and run a program, and wait for it \\
\bottomrule
\end{tabularx}
\end{center}

\section{Reserved descriptors}

\begin{center}
\begin{tabularx}{0.85\textwidth}{@{}l l X@{}}
\toprule
\textbf{fd} & \textbf{Name} & \textbf{Device} \\
\midrule
0 & stdin & Keyboard (blocks until Enter) \\
1 & stdout & Screen \\
2 & stderr & Screen \\
$\geq 3$ & --- & Open files \\
\bottomrule
\end{tabularx}
\end{center}

% ============================================================
\chapter{Source file index}

\begin{center}
\begin{tabularx}{\textwidth}{@{}p{5.0cm} p{1.2cm} X@{}}
\toprule
\textbf{File} & \textbf{Ch.} & \textbf{Responsibility} \\
\midrule
\cod{boot/bootloader.asm} & 12--15 & Boot sector: print, E820, load the kernel, jump \\
\cod{kernel/kernel\_intro.asm} & 16--18 & GDT, protected mode, kernel stack, enter C \\
\cod{kernel/linker.ld} & 19 & Places the kernel at \hexa{9000} \\
\cod{kernel/kernel\_main.c} & 19 & Subsystem initialisation and self-tests \\
\cod{kernel/kernel\_utils.c} & 19 & \cod{memset}, \cod{memcpy}, \cod{strcmp}, \cod{strncpy} \\
\cod{kernel/interrupts.asm} & 21 & ISR/IRQ stubs, port I/O, syscall entry \\
\cod{kernel/interrupts.c} & 21--22 & IDT setup, PIC remap, exception and syscall dispatch \\
\cod{kernel/drivers/vga\_text.c} & 19--20 & Text mode output \\
\cod{kernel/drivers/timer.c} & 23 & PIT, tick counter, wake sleepers, call the scheduler \\
\cod{kernel/drivers/keyboard.c} & 24 & Scancodes, modifiers, delivery to blocked readers \\
\cod{kernel/drivers/ata.c} & 30 & Disk read and write, PIO + LBA \\
\cod{memory/pmm.c} & 25 & Frame bitmap: allocate, free, protect \\
\cod{memory/vmm.c} & 26--28 & Paging: map, unmap, translate, fault handler \\
\cod{memory/vmm.asm} & 26--27 & \reg{CR0}, \reg{CR2}, \reg{CR3}, TLB \\
\cod{memory/heap.c} & 29 & \cod{kmalloc} / \cod{kfree} \\
\cod{tasks/procman.c} & 30--32 & Process creation, destruction, TSS \\
\cod{tasks/procman.asm} & 32 & \cod{ltr} \\
\cod{tasks/context.c} & 31 & Save and restore context, switch \\
\cod{tasks/scheduler.c} & 31 & Round robin, quantum \\
\cod{tasks/loader.c} & 33 & Read a binary and map it into a new process \\
\cod{filesystem/fs.c} & 31--32 & Blocks, inodes, directories, read and write \\
\cod{filesystem/fs\_manager.c} & 32 & Path resolution, descriptors, \cod{ls}/\cod{mkdir}/\cod{rm} \\
\cod{tools/mkfs.c} & 36 & Host-side formatter and \cod{rootfs} installer \\
\cod{user/libc/syscalls.asm} & 34 & The \cod{int 0x80} bridge \\
\cod{user/libc/output.c} & 35 & \cod{printf} and friends \\
\cod{user/libc/memory.c} & 35 & \cod{malloc} on top of \cod{sbrk} \\
\cod{user/programs/shell.c} & 37 & The command interpreter \\
\bottomrule
\end{tabularx}
\end{center}

% ============================================================
\chapter{Notes on the exercises}

These are not full solutions --- they are the shape of the answer, so you can check whether your
reasoning went the right way.

\begin{description}[leftmargin=0pt,style=nextline,font=\bfseries\color{accent}]

\item[1.1 Why 512 bytes cannot hold an OS]
512 bytes is roughly 150 machine instructions. Enough to print a message and read from disk, and
nothing else. The minimum it must do is: locate more code on the disk, load it into memory, and jump
to it. Everything else is optional.

\item[2.1 Why divide by 32 and take the remainder]
The bitmap is an array of 32-bit words. \cod{i / 32} selects the word holding frame $i$; \cod{i \% 32}
selects the bit within that word. Together they turn a flat bit index into a (word, bit) pair.

\item[3.1 Why CR3 holds a physical address]
If it held a virtual address, translating it would require reading the page directory --- which you
cannot find without translating \reg{CR3} first. Infinite regress. The chain has to start somewhere
untranslated, and \reg{CR3} is that starting point.

\item[4.1 Why the boot sector loads at \hexa{7C00}]
Below \hexa{500} is the interrupt vector table and BIOS data. \hexa{7C00} is $32\,\text{KB} - 1024$:
as high as possible in the 32 KB of the original PC while leaving the last kilobyte free. It also
leaves a usable gap above it, which is where this project puts the kernel.

\item[12.1 Why \cod{print\_string} uses SI]
\cod{lodsb} implicitly reads from \cod{DS:SI} and increments \cod{SI}. It only works with that
register. Using \cod{BX} would require a separate load and an explicit increment.

\item[16.1 \cod{cs = 0x18} without the RPL]
The processor checks that the RPL is no more privileged than the descriptor's DPL, and \cod{iret}
will not lower privilege unless the target \reg{CS} has RPL 3. With \cod{cs = 0x18} (RPL 0) pointing
at a DPL 3 descriptor, the \cod{iret} raises a general protection fault instead of entering ring 3.

\item[25.1 Is the E820 buffer at \hexa{5000} a bug?]
No, though it looks like one. The region is marked usable and never explicitly protected, so
\cod{alloc\_frame} can hand it out. But the map is read exactly once, inside \cod{init\_pmm}, which
runs before any allocation happens. After that the data is dead. It is safe by ordering rather than
by design --- which is worth a comment in the source.

\item[26.1 Splitting \hexa{C0000000}]
$\hexa{C0000000} \gg 22 = 768$. Table index and offset are both 0. So the kernel heap starts exactly
at the beginning of directory entry 768, meaning it occupies whole 4 MB directory slots from there
up. That alignment is why the constant is \hexa{C0000000} and not some nearby value: it keeps the
kernel heap from sharing a directory entry with anything else.

\item[7.1 Which calls \cod{cat file.txt} needs]
\cod{open} (\hexa{09}) to get a descriptor, then \cod{read} (\hexa{06}) in a loop, writing each chunk
with \cod{write} (\hexa{05}) to descriptor 1, then \cod{close} (\hexa{0A}). Four of the ten system
calls, which is a good illustration of how little a useful program actually needs.

\item[20.1 \cod{vga\_text\_write\_hex}]
Loop from the most significant nibble down: shift right by $4i$ for $i = 7 \ldots 0$, mask with
\hexa{F}, and convert with \cod{d < 10 ? '0' + d : 'A' + d - 10}. Print \cod{"0x"} first. The
difference from decimal is that hexadecimal has a fixed width, so you print leading zeros rather than
suppressing them --- which also means no reversal buffer is needed.

\end{description}

% ============================================================
\chapter{The runnable examples}

Every recipe box in this book points at a folder under
\filepath{doc/claudedoc/examples/}. Each folder holds a small self-contained program plus a
Makefile that builds a complete bootable image and starts QEMU.

\begin{center}
\begin{tabularx}{\textwidth}{@{}l l X@{}}
\toprule
\textbf{Folder} & \textbf{Ch.} & \textbf{Demonstrates} \\
\midrule
\cod{ch12-bootloader}      & 12 & Editing the 512-byte boot sector and printing via the BIOS \\
\cod{ch16-gdt}             & 16 & Reading the live GDT with \cod{SGDT}; the Accessed and Busy bits \\
\cod{ch18-vga-output}      & 18 & Every text-driver operation, and the \cod{write\_dec(0)} bug \\
\cod{ch21-custom-interrupt}& 21 & Adding your own IDT vector, in all three layers \\
\cod{ch23-timer}           & 23 & Timing work, and what changing the PIT frequency really changes \\
\cod{ch24-keyboard-input}  & 24 & Reading a line correctly, unlike the shell \\
\cod{ch25-physical-memory} & 25 & \cod{alloc\_frame} / \cod{free\_frame} and first-fit reuse \\
\cod{ch27-paging}          & 27 & Map, alias, translate, unmap, and build an address space \\
\cod{ch28-page-fault}      & 28 & Touch an unmapped address and read the report \\
\cod{ch29-heap}            & 29 & \cod{kmalloc}/\cod{kfree} with the block list printed each step \\
\cod{ch30-kernel-process}  & 30 & Three kernel processes preempted by the timer \\
\cod{ch33-loader}          & 33 & Loading a program from the kernel, with no shell \\
\cod{ch34-new-syscall}     & 34 & Adding a system call: the four files you must touch \\
\cod{ch35-libc}            & 35 & A user-program template using the whole libc \\
\cod{ch37-ata-sectors}     & 37 & Raw sector read and write, and finding the superblock \\
\cod{ch39-filesystem}      & 39 & Create, write, reopen, read, EOF, list, delete \\
\cod{ch40-shell-command}   & 40 & Adding \cod{echo} and \cod{rev} to the shell \\
\cod{ch45-reboot-loop}     & 45 & Reproducing the reboot loop, and tracing it \\
\bottomrule
\end{tabularx}
\end{center}

\section{Running one}

\begin{lstlisting}[style=sh]
cd doc/claudedoc/examples/ch27-paging
make run      # build and boot it in a QEMU window
make shell    # boot headless and print the text screen (works over ssh)
make trace    # boot with interrupt logging, stop at a triple fault
make diff     # show exactly what this example changes
make clean    # delete this example's build directory
\end{lstlisting}

\cod{make trace} is the one to reach for when something crashes: it is the technique from
Chapter 42 wired into every example.

\begin{concept}{The project is never modified}
Each example copies the OS tree into its own \filepath{build/os/}, overlays its files on the
copy, builds that and boots it. \filepath{lytlnybl/} is left untouched, and deleting
\filepath{build/} leaves nothing behind. You can break an example as thoroughly as you like.
\end{concept}

\section{How an example changes the system}

Two mechanisms, chosen per example.

\textbf{Overlay} replaces a whole file, and is what most examples use --- nearly always
\cod{kernel\_main.c}, since that is where the kernel's own demonstration code lives:

\begin{lstlisting}[style=mk]
DESC    = Chapter 27 -- mapping, aliasing, translating and unmapping pages
OVERLAY = kernel_main.c:kernel/kernel_main.c

include ../common/common.mk
\end{lstlisting}

\textbf{Patches} insert into existing files, for changes that touch several places at once:

\begin{lstlisting}[style=mk]
DESC        = Chapter 34 -- adding a system call, end to end (4 files)
OVERLAY     = test.c:user/programs/test.c
PATCHES     = patch.py
SCREEN_ARGS = --wait 10 --type "run /bin/user_test\n" --settle 4

include ../common/common.mk
\end{lstlisting}

Every patch is anchored on exact existing text and \emph{fails loudly} if that text is missing or
appears more than once, so an example can never seem to apply and then silently do nothing. Read
\filepath{ch34-new-syscall/patch.py} as the recipe itself: its four edits are the four edits
adding a system call always requires.

\section{Checking they still build}

\begin{lstlisting}[style=sh]
cd doc/claudedoc/examples && ./verify.sh
\end{lstlisting}

Builds all eighteen from scratch and reports. Worth running after changing anything in the OS an
example overlays, since an overlay can go stale if the file it replaces gains something the example
does not know about.

\begin{warn}{What \cod{make shell} needs}
It boots QEMU with no display, asks the monitor for a screenshot, and matches each 9$\times$16 cell
against the VGA font it finds inside the vgabios ROM that ships with QEMU --- recovering the screen
as text. It therefore needs Python 3 and that ROM
(\filepath{/usr/share/seabios/vgabios.bin} on Debian and Ubuntu). If it cannot find one it says
so, and \cod{make run} works regardless.
\end{warn}
